\section{Results}

\subsection{Derived FOPDT transfer functions}
The transfer functions of the system were derived using two different methods, namely the process reaction curve method and transfer function fitting using MATLAB. Each method
yielded a MIMO transfer function matrix with the following mathematical description:

\begin{equation}
    \begin{bmatrix}
        Y_1(s) \\
        Y_2(s)
    \end{bmatrix}
    =
    \begin{bmatrix}
        G_{11}(s) & G_{12}(s) \\
        G_{21}(s) & G_{22}(s)
    \end{bmatrix}
    \begin{bmatrix}
        U_1(s) \\
        U_2(s)
    \end{bmatrix}
\end{equation}

The following table contains these two matrices.

\begin{table}[htbp]
    \centering
    % Increases the vertical space so the fractions don't squash together
    \renewcommand{\arraystretch}{2.5} 
    \caption{Comparison of Transfer Functions: Transfer Function Fitting vs. Process Reaction Curve}
    \label{tab:transfer_function_comparison}
    \begin{tabular}{lcc}
        \hline
        \textbf{Thermal Pathway} & \textbf{Transfer Function Fitting (\texttt{tfest})} & \textbf{Process Reaction Curve} \\
        \hline
        
        % Row 1: G11 (PWM1 to T1)
        $G_{11}(s)$ (PWM 1 $\rightarrow$ T1) & 
        $\frac{-0.1974s + 0.05300}{s^2 + 0.4582s + 0.0034}$ & 
        $\frac{K_{11}}{\tau_{11} s + 1} e^{-\theta_{11} s}$ \\
        
        % Row 2: G12 (PWM2 to T1)
        $G_{12}(s)$ (PWM 2 $\rightarrow$ T1) & 
        $\frac{-0.0156 s + 0.0010}{s^2 + 0.0300 s + 0.0002}$ & 
        $\frac{K_{12}}{\tau_{12} s + 1} e^{-\theta_{12} s}$ \\
        
        % Row 3: G21 (PWM1 to T2)
        $G_{21}(s)$ (PWM 1 $\rightarrow$ T2) & 
        $\frac{-0.0097 s + 0.0016}{s^2 + 0.0268 s + 0.0002}$ & 
        $\frac{K_{21}}{\tau_{21} s + 1} e^{-\theta_{21} s}$ \\
        
        % Row 4: G22 (PWM2 to T2)
        $G_{22}(s)$ (PWM 2 $\rightarrow$ T2) & 
        $\frac{-0.0346 s + 0.0234}{s^2 + 0.3998 s + 0.0024}$ & 
        $\frac{K_{22}}{\tau_{22} s + 1} e^{-\theta_{22} s}$ \\
        
        \hline
    \end{tabular}
\end{table}


\textbf{Comment on the table}

\subsection{Matlab fitting}
The \textbf{tfest} function was used in Matlab to fit the data to a transfer function. The following steps were followed to achieve this:
\\
\\
\textbf{Data Preparation:}
\\
The recorded data was sliced to include a sufficient amount of zero inputs, this is required by the tfest function to properly fit the model. After which the data was broken into a voltage vector (the PWM input), and two temperature vectors 
(temperatures of each of the transistors). The ambient temperature was subtracted from  the temperature vectors and the data was packaged into an iddata object.
\\
\\
\textbf{Model Fitting:}
\\
Three parameters were passed to the tfest function, namely: the iddata object, the number of poles the system has, and the number of zeros the system has.
The number of poles a system has is determined by the number of energy storing elements, in the case of this system the transistors are storing energy in the form
of heat, thus the system has two poles. It is common practise to take the number of zeros of an unknown system as: Zeros $=$ Poles $- 1$. 
The fitted models were then plotted using the compare function as can be seen in the Figure \ref{fig:pwm1_comparison} and Figure \ref{pwm2_comparison} below.



\begin{figure}[htbp]
    \centering
    \includegraphics[width=1\textwidth]{Images/PWM1Step_Compare.pdf} 
    \caption{Comparison of the physical sensor data against the simulated SIMO transfer function model for PWM 1.}
    \label{fig:pwm1_comparison}
\end{figure}

Figure \ref{fig:pwm1_comparison} displays the fitted model for both $G_{11}$ (top) and $G_{21}$ (bottom) which are similarly accurate; 94.1\% and 94.4\%, on top
of the measured physical data.

\begin{figure}[htbp]
    \centering
    \includegraphics[width=1\textwidth]{Images/PWM2Step_Compare.pdf} 
    \caption{Comparison of the physical sensor data against the simulated SIMO transfer function model for PWM 2.}
    \label{fig:pwm2_comparison}
\end{figure}



\begin{figure}[htbp]
    \centering
    \includegraphics[width=1\textwidth]{Images/PWM1Step_TestData.pdf} 
    \caption{Comparison of the physical sensor data for a 50\% duty cycle test case against the simulated SIMO transfer function model for PWM 1.}
    \label{fig:pwm1_test}
\end{figure}

\begin{figure}[htbp]
    \centering
    \includegraphics[width=1\textwidth]{Images/PWM2Step_TestData.pdf} 
    \caption{Comparison of the physical sensor data for a 50\% duty cycle test case against the simulated SIMO transfer function model for PWM 2.}
    \label{fig:pwm2_test}
\end{figure}